\documentclass[11pt,letterpaper]{article}
\usepackage[T1]{fontenc}
\usepackage{lmodern,microtype}
\usepackage[margin=.75in,headheight=15pt]{geometry}
\usepackage{amsmath,amssymb,mathtools}
\usepackage{xcolor,enumitem,fancyhdr,hyperref}
\definecolor{ink}{HTML}{222834}
\definecolor{indigo}{HTML}{25355D}
\definecolor{muted}{HTML}{505A6B}
\definecolor{keyink}{HTML}{205D52}
\color{ink}
\hypersetup{colorlinks=true,linkcolor=indigo,urlcolor=indigo,
  pdftitle={Probability Midterm 1},pdfauthor={Igor Rivin}}
\setlength{\parindent}{0pt}
\setlength{\parskip}{4pt}
\setlength{\abovedisplayskip}{6pt}
\setlength{\belowdisplayskip}{6pt}
\clubpenalty=10000
\widowpenalty=10000
\setlist[enumerate]{label=(\alph*),leftmargin=1.7em,
  itemsep=4pt,topsep=5pt,parsep=2pt}
\newcommand{\Prob}{\mathbb P}
\newcommand{\E}{\mathbb E}
\DeclareMathOperator{\Var}{Var}
\newcommand{\edition}{\ifanswers Instructor key\else Student paper\fi}
\newcommand{\partpts}[1]{\unskip\nobreak\hspace{.3em}%
  \mbox{\small\color{muted}(#1 points)}}
\newcommand{\credit}[1]{\par{\footnotesize\color{muted}\textit{Credit:} #1}}
\newcommand{\problem}[3]{\par\addvspace{9pt}%
  {\large\bfseries\color{indigo}#1.\ #2}\hfill
  {\small\color{muted}\textbf{#3 points}}\par\nobreak}
\NewDocumentEnvironment{solution}{O{1in}+b}{%
  \ifanswers
    \par\begingroup\small\color{keyink}\textbf{Solution.} #2\par\endgroup
  \else
    \par\nobreak\vspace*{#1}
  \fi
}{}
\pagestyle{fancy}
\fancyhf{}
\fancyhead[L]{\small\color{muted}Probability \textbar\ Midterm 1}
\fancyhead[R]{\small\color{muted}Chance, Data, and Decisions}
\fancyfoot[L]{\footnotesize\color{muted}\edition}
\fancyfoot[R]{\footnotesize\color{muted}Page \thepage\ of \pageref*{exam-last-page}}
\renewcommand{\headrulewidth}{.3pt}
\begin{document}
{\LARGE\bfseries\color{indigo}Probability Midterm 1}\par
{\small\color{muted}50 minutes \quad $\bullet$ \quad 100 points
  \quad $\bullet$ \quad Six problems}\par
\ifanswers
\textbf{Grading note.} Award method credit after an earlier error; do not
deduct twice for the same mistake. Suggested working times for Problems
1--6 are 5, 5, 6, 8, 10, and 9 minutes, leaving 7 minutes to check.
\else
\smallskip
\textbf{Name:}\ \rule{.49\linewidth}{.4pt}\hfill
\textbf{Date:}\ \rule{.20\linewidth}{.4pt}\par
\fi
\textbf{Instructions.} Answer all six problems and show your reasoning.
Exact answers are acceptable; binomial coefficients need not be evaluated.
No calculator is needed. Use the spaces provided; clearly label any work
continued on the back. In Problem 6, use the guided steps rather than
quoting a coupon-collector formula. Reserve time to check your work.

\problem{1}{Sets}{10}
Let $U=\{1,2,\ldots,10\}$, $A=\{2,4,6,8,10\}$, and
$B=\{1,2,3,4\}$. Take all complements relative to $U$.
\begin{enumerate}
\item Find $A\cap B$, $A\setminus B$, and $(A\cup B)^c$.
\partpts{6}
\begin{solution}[.85in]
\[
A\cap B=\{2,4\},\qquad A\setminus B=\{6,8,10\},\qquad
(A\cup B)^c=\{5,7,9\}.
\]
\credit{2 points per correct set; allow partial credit for a single entry error.}
\end{solution}
\item For these sets $A,B$, is the following statement true for every
$C\subseteq U$? Justify your answer with a proof or a counterexample.
\partpts{4}
\[
C\subseteq A\cup B\quad\Longrightarrow\quad
C\subseteq A\ \text{or}\ C\subseteq B.
\]
\begin{solution}[.85in]
False. Take $C=\{1,6\}$. Both elements lie in $A\cup B$, so
$C\subseteq A\cup B$. But $1\notin A$ and $6\notin B$;
therefore $C\nsubseteq A$ and $C\nsubseteq B$.
\credit{False: 1; valid counterexample: 1; verifying containment in the union: 1;
verifying both failed containments: 1. Accept any valid counterexample.}
\end{solution}
\end{enumerate}

\problem{2}{A probability space}{10}
Let $\Omega=\{a,b,c,d\}$, with every subset of $\Omega$ an event.
The probabilities of the individual outcomes $a,b,c,d$ are
$k,2k,3k,4k$, respectively.
\begin{enumerate}
\item Find $k$, explaining the probability axiom you use.
\partpts{4}
\begin{solution}[.6in]
The disjoint singleton events partition $\Omega$, and $\Prob(\Omega)=1$.
Finite additivity gives $k+2k+3k+4k=1$, so $k=1/10$.
All four probabilities are nonnegative.
\credit{Normalization/additivity explanation: 2; equation: 1; $k=1/10$: 1.}
\end{solution}
\item Let $E=\{a,c\}$ and $F=\{b,c\}$.
Find $\Prob(E)$, $\Prob(E\cap F)$, and $\Prob(E^c)$.
\partpts{6}
\begin{solution}[.85in]
Since $E\cap F=\{c\}$ and $E^c=\{b,d\}$,
\[
\Prob(E)=\frac{1+3}{10}=\frac25,\qquad
\Prob(E\cap F)=\frac3{10},\qquad
\Prob(E^c)=\frac{2+4}{10}=\frac35.
\]
\credit{2 points per probability. If $k$ was incorrect in (a), award full
method credit for $4k,3k,6k$, or a consistently justified complement.}
\end{solution}
\end{enumerate}

\newpage
\problem{3}{Inclusion--exclusion}{15}
In a class of 40 students, 18 study French, 16 Spanish, and 12 German.
Seven study both French and Spanish, five both French and German,
and four both Spanish and German. Two study all three languages.
\textbf{The pairwise counts include students who study all three.}
\begin{enumerate}
\item How many students study at least one of these languages?
Show your inclusion--exclusion calculation. \partpts{11}
\begin{solution}[.95in]
Writing $F,S,G$ for the three sets, inclusion--exclusion gives
\begin{align*}
|F\cup S\cup G|
&=|F|+|S|+|G|-|F\cap S|-|F\cap G|-|S\cap G|
  +|F\cap S\cap G|\\
&=18+16+12-7-5-4+2=32.
\end{align*}
\credit{Adding the three individual counts: 3; subtracting each pairwise
intersection: 3; restoring the triple intersection once: 3; result 32: 2.}
\end{solution}
\item How many students study none of these languages? \partpts{4}
\begin{solution}[.65in]
There are $40-|F\cup S\cup G|=40-32=8$ such students.
\credit{Complement subtraction: 3; answer 8: 1. Carry forward an earlier error.}
\end{solution}
\end{enumerate}

\problem{4}{Balls in an urn}{20}
An urn contains five red, four blue, and three green balls. Four balls
are selected uniformly without replacement: every four-ball subset
of the twelve individual balls is equally likely.
\begin{enumerate}
\item Find the probability of selecting exactly two red, one blue,
and one green ball. \partpts{10}
\begin{solution}[1.25in]
There are $\binom{12}{4}$ equally likely subsets. Choose two of the
five red balls, one of the four blue balls, and one of the three green balls:
\[
\Prob(\text{two red, one blue, one green})
=\frac{\binom52\binom41\binom31}{\binom{12}4}
=\frac{120}{495}=\frac8{33}.
\]
\credit{Uniform denominator: 3; red choice: 2; blue and green choices: 2;
multiplying the counts and forming the ratio: 3. No numerical simplification needed.}
\end{solution}
\item Find the probability of selecting at least one green ball.
\partpts{10}
\begin{solution}[1.25in]
There are nine non-green balls. The complementary event selects all
four balls from these nine, so
\[
\Prob(\text{at least one green})
=1-\frac{\binom94}{\binom{12}4}
=1-\frac{126}{495}=\frac{41}{55}.
\]
\credit{Complement event: 3; count $\binom94$: 3; denominator: 2;
subtracting the ratio from 1: 2. Accept an equivalent direct count.}
\end{solution}
\end{enumerate}

\newpage
\problem{5}{A density, its mean, and its variance}{25}
A continuous random variable $X$ has probability density
\[
f(x)=\begin{cases}
cx^3,&0\le x\le1,\\
0,&\text{otherwise},
\end{cases}
\]
where $c$ is a constant. Show the integrals you use.
\begin{enumerate}
\item Find $c$. \partpts{5}
\begin{solution}[1.4in]
A density must integrate to 1:
\[
1=\int_{-\infty}^{\infty}f(x)\,dx
=c\int_0^1 x^3\,dx
=c\left[\frac{x^4}{4}\right]_0^1=\frac c4.
\]
Thus $c=4$, which also makes $f$ nonnegative.
\credit{Normalization integral: 2; integral value: 2; $c=4$: 1.}
\end{solution}
\item Compute $\E[X]$. \partpts{8}
\begin{solution}[1.5in]
Using $f(x)=4x^3$ on $[0,1]$,
\[
\E[X]=\int_{-\infty}^{\infty}x f(x)\,dx
=4\int_0^1 x^4\,dx
=4\left[\frac{x^5}{5}\right]_0^1=\frac45.
\]
\credit{Moment integral with the extra factor $x$: 4;
antiderivative and limits: 3; result: 1. Carry forward an earlier value of $c$.}
\end{solution}
\item Compute $\E[X^2]$, and hence $\Var(X)$. \partpts{12}
\begin{solution}[2in]
\[
\E[X^2]=\int_{-\infty}^{\infty}x^2 f(x)\,dx
=4\int_0^1 x^5\,dx
=4\left[\frac{x^6}{6}\right]_0^1=\frac23.
\]
Consequently,
\[
\Var(X)=\E[X^2]-(\E[X])^2
=\frac23-\left(\frac45\right)^2
=\frac{50-48}{75}=\frac2{75}.
\]
\credit{Second-moment integral: 4; evaluating it: 3;
variance identity: 3; substitution and arithmetic: 2.
Accept a correctly evaluated direct variance integral as an alternative.}
\end{solution}
\end{enumerate}

\newpage
\problem{6}{Collecting four toy types}{20}
Each cereal box contains one toy. There are four toy types, each equally
likely, and different boxes have independent contents. Duplicates are possible.
Starting with no toys, let $T$ be the number of boxes opened up to and
including the first box that completes the collection of all four types.
\begin{enumerate}
\item After exactly $k$ distinct types have been collected, let $p_k$
be the probability that the next box contains a new type. Find
$p_0,p_1,p_2,p_3$, and explain your reasoning. \partpts{6}
\begin{solution}[1.3in]
Exactly $4-k$ of the four types are new, so $p_k=(4-k)/4$.
Thus
\[
p_0=1,\qquad p_1=\frac34,\qquad p_2=\frac12,\qquad p_3=\frac14.
\]
\credit{Counting the missing types and dividing by four: 2;
four correct probabilities: 1 each.}
\end{solution}
\item Let $W_k$ be the number of additional boxes needed to go from
exactly $k$ collected types to $k+1$ collected types, including the
box containing the new type. Find $\E[W_k]$ for $k=0,1,2,3$.
\partpts{8}

\textit{You may use:} A geometric waiting time that counts trials
through the first success, with success probability $p$, has mean $1/p$.
\begin{solution}[1.6in]
Until a new type appears, the set of missing types stays the same.
Independence of boxes makes this a geometric wait with success
probability $p_k$. For $k=0$, the first box always supplies a new type.
Therefore
\[
\E[W_0]=1,\qquad \E[W_1]=\frac43,\qquad
\E[W_2]=2,\qquad \E[W_3]=4.
\]
\credit{2 points per correctly obtained mean. Carry forward earlier
probabilities when the reciprocal rule is applied consistently.}
\end{solution}
\item Express $T$ in terms of $W_0,W_1,W_2,W_3$ and compute $\E[T]$.
State the expectation rule you use. \partpts{6}
\begin{solution}[1.8in]
The four successive stages partition all boxes opened, so
$T=W_0+W_1+W_2+W_3$. Linearity of expectation gives
\[
\E[T]=\sum_{k=0}^3\E[W_k]
=1+\frac43+2+4=\frac{25}{3}.
\]
This is an expected number of boxes, not a probability, so it need not
be an integer. Independence of the $W_k$ is not needed to apply linearity.
\credit{Stage decomposition: 2; linearity and the sum of stage means: 2;
evaluation: 2. Do not require a proof of independence of the waits.}
\end{solution}
\end{enumerate}
\vfill
\begin{center}\small\color{muted}End of paper. Check that you have answered all six problems.\end{center}
\label{exam-last-page}
\end{document}
